% TGL-Rec: Time-Aware Graph Evidence for LLM-Based Recommendation
% Target: Top-conference (NeurIPS/ICML/ICLR/WWW/KDD)

\documentclass[10pt,twocolumn]{article}
\usepackage[utf8]{inputenc}
\usepackage{amsmath,amssymb}
\usepackage{booktabs}
\usepackage{graphicx}
\usepackage{hyperref}
\usepackage{algorithm}
\usepackage{algorithmic}

\title{Do LLM Recommenders Follow Time or Similarity?\\
Temporal Graph Evidence for Need-Aware Sequential Recommendation}

\author{Anonymous}
\date{}

\begin{document}
\maketitle

\begin{abstract}
Large language model (LLM)-based recommenders serialize user interaction histories into prompts, yet our diagnostic experiments reveal they are largely insensitive to temporal order: shuffling or reversing histories produces near-identical rankings. This exposes a \emph{semantic trap}---the tendency to recommend items similar to recent interactions rather than items matching the user's evolving temporal need. We propose \textbf{TGL-Rec}, a framework that constructs a Temporal Directed Item Graph (TDIG) from train-only interactions, learns a lightweight \emph{need-gate} to distinguish temporal-transition signals from semantic-similarity signals, translates graph evidence into auditable natural language, and fine-tunes a local LLM reranker via LoRA. On four Amazon domains under a shared 101-candidate protocol with 8 official LLM4Rec baselines, TGL-Rec achieves state-of-the-art performance, with the need-gate and temporal evidence each contributing measurable gains confirmed by ablation.
\end{abstract}

\section{Introduction}

Sequential recommendation assumes that the order and timing of user interactions carry predictive signal for future behavior. Recent LLM-based recommenders~\cite{llm2rec,llmesr,rlmrec} leverage large language models to process interaction histories as text, achieving strong results on standard benchmarks. However, a fundamental question remains underexplored: \emph{do these models actually use temporal order, or do they primarily rely on semantic similarity between item descriptions?}

We investigate this question through controlled perturbation experiments. Specifically, we shuffle, reverse, and truncate user histories before prompting Qwen3-8B for candidate ranking. Our finding is striking: across four Amazon product domains, shuffling the interaction history changes the model's MRR by less than 2\%, suggesting that the base LLM treats the history as a bag of items rather than a temporal sequence.

This observation motivates our core contribution: if LLMs ignore temporal order, we must \emph{explicitly provide} temporal evidence rather than hoping the model will extract it from serialized histories. We introduce three mechanisms:

\begin{enumerate}
\item \textbf{Temporal Directed Item Graph (TDIG)}: A train-only directed graph where edges represent consecutive user transitions, annotated with time-gap statistics, transition probabilities, and directional asymmetry.

\item \textbf{Need-Aware Gate}: A lightweight learned gate ($\sim$25 parameters) that predicts, for each user-candidate pair, whether temporal-transition evidence or semantic-similarity evidence should dominate the ranking decision. This detects \emph{semantic traps}---candidates that appear similar to recent items but do not match the user's temporal need pattern.

\item \textbf{Graph-to-Language Evidence}: Deterministic templates that translate TDIG paths into natural-language evidence statements, grounding the LLM reranker's decisions in auditable graph statistics rather than implicit pattern matching.
\end{enumerate}

Unlike prior graph-to-language approaches (e.g., G-Refer) that use static collaborative filtering paths, TGL-Rec operates on \emph{temporal} directed graphs with time-decay, windowed co-occurrence, and need-state encoding. Unlike LLM enhancement methods (e.g., LLM-ESR) that augment sequential models, TGL-Rec provides \emph{explicit evidence} to the LLM about when and why temporal transitions matter.

We evaluate TGL-Rec against 8 official LLM4Rec baselines on four Amazon domains (Beauty, Books, Electronics, Movies) under a shared same-candidate protocol with 101 candidates per user. All methods use Qwen3-8B as the backbone. Our ablation study confirms that each component---the temporal graph, the need-gate, the semantic trap penalty, and the evidence translation---contributes to the final performance.

\section{Related Work}

\paragraph{LLM-Based Recommendation.}
% LLM2Rec, LLM-ESR, LLMEmb, RLMRec, IRLLRec, ELMRec, ProEx, ProMax
% Key point: none explicitly model temporal transitions as evidence

\paragraph{Sequential and Time-Aware Recommendation.}
% SASRec, BERT4Rec, TiSASRec, time-aware models
% Key point: these use temporal signals but not as explicit LLM evidence

\paragraph{Graph-Based Recommendation.}
% G-Refer, LightGCN, knowledge graph approaches
% Key point: G-Refer uses static paths, we use temporal directed graphs

\paragraph{Diagnostic Evaluation of Recommenders.}
% Lost in Sequence, perturbation studies
% Key point: we extend diagnostics to temporal sensitivity

\section{Observation: LLMs Ignore Temporal Order}

% Present perturbation experiment results here
% Table: base vs shuffled vs reversed vs recent_only
% Key finding: MRR difference < 2% between base and shuffled

\section{Method: TGL-Rec}

\subsection{Temporal Directed Item Graph}
% Construction from train-only interactions
% Edge statistics: transition count, probability, PMI, lift, direction asymmetry
% Time-window and time-decay annotations

\subsection{Need-Aware Gate}
% Input: user need-state vector + candidate evidence features
% Output: alpha in [0,1] (temporal vs semantic balance)
% Training: positive = ground truth with temporal evidence, negative = candidates without
% Semantic trap detection: high semantic similarity + low temporal support

\subsection{Graph-to-Language Evidence Translation}
% Deterministic templates (no hallucination)
% Evidence types: transition, time_window, semantic, contrastive, history, user_drift
% Token budget and diversity constraints

\subsection{LoRA Reranking with Evidence}
% Qwen3-8B + QLoRA (r=16, alpha=32)
% Training data: top-K candidates scored by gate, evidence text included
% Inference: same chat template, evidence-grounded ranking

\section{Experiments}

\subsection{Setup}
% 4 Amazon domains, same-candidate protocol (101 candidates)
% 8 official baselines (all Qwen3-8B backbone)
% Metrics: MRR, HR@5, HR@10, NDCG@5, NDCG@10
% 5 seeds for significance

\subsection{Main Results}
% Table placeholder

\subsection{Ablation Study}
% 7 ablation variants
% Table placeholder

\subsection{Analysis}
% When does temporal evidence help most?
% Semantic trap detection accuracy
% Gate behavior visualization

\section{Conclusion}

% Summary of findings and contributions

\end{document}
